\chapter{串联电路}\label{ch:g4:series-circuits}

去年，一串彩灯因为缺了一只
\omterm{def:g3:first-electric-circuit:bulb}{灯泡}而全部熄灭——我们当时承诺过，这种连接方式有个名字。
名字就是\omterm{def:g4:series-circuits:series}{串联电路}：把许多用电器接在同一个回路上，这是最简单的办法。
它有严格的规则，你抽屉里的手电筒
一直都在遵守它们。

\section{一个回路，多个用电器}

\begin{definition}[串联]\label{def:g4:series-circuits:series}
电路中的用电器若\emph{串联}\index{串联电路}，
即它们一个接一个地排列在同一个回路上，
就像一串项链上的珠子：电流的路径依次经过第一个、第二个、第三个，
中间没有任何分岔。按这种方式连接的电路，
称为\emph{串联电路}。
\end{definition}

\begin{omfigure}
\begin{tikzpicture}[scale=0.9]
  \draw (0,0) to[battery1, l=电池] (0,2.6)
    to[short] (1.6,2.6)
    to[cspst, l=开关] (3.2,2.6)
    to[lamp, l=灯泡 1] (5.2,2.6)
    to[lamp, l=灯泡 2] (7.2,2.6)
    to[short] (7.2,0)
    to[short] (0,0);
\end{tikzpicture}

{\small \omterm{def:g4:series-circuits:series}{串联电路}：\omterm{def:g3:first-electric-circuit:battery}{电池}、\omterm{ex:g3:first-electric-circuit:switch}{开关}和两只\omterm{def:g3:first-electric-circuit:bulb}{灯泡}位于同一个回路上——
每个用电器都穿在同一根项链上。}
\end{omfigure}

\begin{method}[搭建电路]\label{met:g4:series-circuits:build}
用一节\omterm{def:g3:first-electric-circuit:battery}{电池}、两只\omterm{def:g3:first-electric-circuit:bulb}{灯泡}、一个\omterm{ex:g3:first-electric-circuit:switch}{开关}和若干导线，按以下步骤连接：
\begin{enumerate}
  \item 用导线把\omterm{def:g3:first-electric-circuit:battery}{电池}的 $+$ \omterm{def:g3:first-electric-circuit:battery}{接线柱}连到
        \omterm{ex:g3:first-electric-circuit:switch}{开关}；
  \item 用导线从\omterm{ex:g3:first-electric-circuit:switch}{开关}连到第一只\omterm{def:g3:first-electric-circuit:bulb}{灯泡}，再从第一只\omterm{def:g3:first-electric-circuit:bulb}{灯泡}
        连到第二只\omterm{def:g3:first-electric-circuit:bulb}{灯泡}；
  \item 闭合回路：从第二只\omterm{def:g3:first-electric-circuit:bulb}{灯泡}连回\omterm{def:g3:first-electric-circuit:battery}{电池}的
        $-$ \omterm{def:g3:first-electric-circuit:battery}{接线柱}；
  \item 合上\omterm{ex:g3:first-electric-circuit:switch}{开关}，观察两只\omterm{def:g3:first-electric-circuit:bulb}{灯泡}是否同时亮起。
\end{enumerate}
开灯前，先用手指沿回路走一遍：如果手指
需要在两条路之间做选择，那就不是\omterm{def:g4:series-circuits:series}{串联电路}。
\end{method}

\section{串联电路的规则}

\begin{proposition}[单一回路的规则]\label{prop:g4:series-circuits:rules}
在\omterm{def:g4:series-circuits:series}{串联电路}中：
\begin{enumerate}
  \item \emph{全有或全无}：任何一处断开——\omterm{ex:g3:first-electric-circuit:switch}{开关}断开，
        \omterm{def:g3:first-electric-circuit:bulb}{灯泡}损坏，导线松脱，所有用电器便
        同时停止工作；
  \item \emph{顺序无关}：交换两只\omterm{def:g3:first-electric-circuit:bulb}{灯泡}，或把
        \omterm{ex:g3:first-electric-circuit:switch}{开关}移到回路的另一侧，一切都与原来
        完全一样；
  \item \emph{分担}：回路上的\omterm{def:g3:first-electric-circuit:bulb}{灯泡}越多，
        每只\omterm{def:g3:first-electric-circuit:bulb}{灯泡}越暗；\omterm{def:g3:first-electric-circuit:battery}{电池}的“推力”
        被整串\omterm{def:g3:first-electric-circuit:bulb}{灯泡}分担。
\end{enumerate}
\end{proposition}

\begin{example}[观察规则]\label{ex:g4:series-circuits:seeing}
搭建双\omterm{def:g3:first-electric-circuit:bulb}{灯泡}电路，逐条检验规则。拧下\emph{任一只}
\omterm{def:g3:first-electric-circuit:bulb}{灯泡}：两只\omterm{def:g3:first-electric-circuit:bulb}{灯泡}都熄灭——这是规则一（被拧下的灯座相当于断点）。
交换两只\omterm{def:g3:first-electric-circuit:bulb}{灯泡}，或把\omterm{ex:g3:first-electric-circuit:switch}{开关}移到二者之间：
亮度毫无变化——这是规则二。现在改用一只\omterm{def:g3:first-electric-circuit:bulb}{灯泡}重新搭建，它
独自发光时比有“伴侣”时明显更亮——
这正是规则三的反向体现。
\end{example}

\begin{example}[彩灯之谜的解答]\label{ex:g4:series-circuits:partylights}
旧日的谜团如今已不再是谜：彩灯上的\omterm{def:g3:first-electric-circuit:bulb}{灯泡}是\omterm{def:g4:series-circuits:series}{串联}的，
几十颗珠子串在同一根项链上。缺一只\omterm{def:g3:first-electric-circuit:bulb}{灯泡}就是一处断点，
规则一使整串彩灯全部熄灭。逐只试\omterm{def:g3:first-electric-circuit:bulb}{灯泡}找出坏的那只，
曾让几代家庭用最笨的办法领会规则一——
现代彩灯串的结构已经不同，
下一年会讲到原因。
\end{example}

\begin{remark}[开关控制所有用电器]\label{rem:g4:series-circuits:switch}
规则二有一个方便的推论：一个\omterm{ex:g3:first-electric-circuit:switch}{开关}，无论放在回路
的哪个位置，都能控制回路上的所有用电器。这就是为什么一个墙上的
\omterm{ex:g3:first-electric-circuit:switch}{开关}可以控制整排的天花板灯——
同样的道理，你无法只关掉\omterm{def:g4:series-circuits:series}{串联}灯串中的\emph{一只}\omterm{def:g3:first-electric-circuit:bulb}{灯泡}，
否则会让全部\omterm{def:g3:first-electric-circuit:bulb}{灯泡}一起熄灭。
\end{remark}

\section{串联更多电池}

\begin{example}[首尾相接]\label{ex:g4:series-circuits:batteries}
\omterm{def:g3:first-electric-circuit:battery}{电池}也可以\omterm{def:g4:series-circuits:series}{串联}排列——不过它们比\omterm{def:g3:first-electric-circuit:bulb}{灯泡}更挑剔：
它们必须\emph{首尾相接}：每节\omterm{def:g3:first-electric-circuit:battery}{电池}的 $+$
接在下一节\omterm{def:g3:first-electric-circuit:battery}{电池}的 $-$ 上。两节\omterm{def:g3:first-electric-circuit:battery}{电池}这样排列时，推力
相加，\omterm{def:g3:first-electric-circuit:bulb}{灯泡}明显比用一节时更亮。把其中一节
\omterm{def:g3:first-electric-circuit:battery}{电池}调个头，变成头对头——两股推力
相互抵消，\omterm{def:g3:first-electric-circuit:bulb}{灯泡}不亮。
\end{example}

\begin{omfigure}
\begin{tikzpicture}[scale=0.8]
  \draw (0,0) to[battery1] (0,1.8)
    to[battery1, l=两节电池] (0,3.6)
    to[short] (3.6,3.6)
    to[lamp, l=更亮！] (3.6,0)
    to[short] (0,0);
  \begin{scope}[xshift=6.9cm]
    \draw (0,0) to[battery1, l=一节电池] (0,1.8)
      to[short] (0,3.6)
      to[short] (3.6,3.6)
      to[lamp, l=更暗] (3.6,0)
      to[short] (0,0);
  \end{scope}
\end{tikzpicture}

{\small 两节首尾相接排成一列的\omterm{def:g3:first-electric-circuit:battery}{电池}共同推动：
同一个\omterm{def:g3:first-electric-circuit:bulb}{灯泡}比只用一节\omterm{def:g3:first-electric-circuit:battery}{电池}时更亮。}
\end{omfigure}

\begin{example}[手电筒内部]\label{ex:g4:series-circuits:flashlight}
打开手电筒，\omterm{def:g3:first-electric-circuit:battery}{电池}一节接一节地滑入
筒中——正好是一列首尾相接的队伍——然后\omterm{def:g3:first-electric-circuit:bulb}{灯泡}和
按钮\omterm{ex:g3:first-electric-circuit:switch}{开关}沿外壳连成一个单一回路。
手电筒\emph{就是}口袋里的\omterm{def:g4:series-circuits:series}{串联电路}。现在你明白了
为什么只要有一节\omterm{def:g3:first-electric-circuit:battery}{电池}装反，手电筒就不亮：
\omterm{def:g3:first-electric-circuit:battery}{电池}头对头，推力相互抵消。
\end{example}

\section{练习}

\begin{exercise}[$\star$]\label{exo:g4:series-circuits:1}
“\omterm{def:g4:series-circuits:series}{串联}”是什么意思？本章借用日常生活中的哪个
情景来说明它？
\end{exercise}

\begin{exercise}[$\star$]\label{exo:g4:series-circuits:2}
说出\omterm{def:g4:series-circuits:series}{串联电路}的三条规则。
\end{exercise}

\begin{exercise}[$\star$]\label{exo:g4:series-circuits:3}
在双\omterm{def:g3:first-electric-circuit:bulb}{灯泡}电路中，\omterm{def:g3:first-electric-circuit:bulb}{灯泡} 2 被拧下，\omterm{def:g3:first-electric-circuit:bulb}{灯泡} 1 会怎样，
为什么？这里起作用的是哪条规则？
\end{exercise}

\begin{exercise}[$\star$]\label{exo:g4:series-circuits:4}
佐伊把\omterm{ex:g3:first-electric-circuit:switch}{开关}从\omterm{def:g3:first-electric-circuit:bulb}{灯泡} 1 前面移到\omterm{def:g3:first-electric-circuit:bulb}{灯泡} 2 后面。
电路的工作情况会发生什么变化？依据哪条规则？
\end{exercise}

\begin{exercise}[$\star$]\label{exo:g4:series-circuits:5}
一节\omterm{def:g3:first-electric-circuit:battery}{电池}能让一个\omterm{def:g3:first-electric-circuit:bulb}{灯泡}正常发光。预测下列情况下的亮度：
两个\omterm{def:g3:first-electric-circuit:bulb}{灯泡}\omterm{def:g4:series-circuits:series}{串联}；\ 三个\omterm{def:g3:first-electric-circuit:bulb}{灯泡}\omterm{def:g4:series-circuits:series}{串联}。被共享的是什么？
\end{exercise}

\begin{exercise}[$\star$]\label{exo:g4:series-circuits:6}
两节\omterm{def:g3:first-electric-circuit:battery}{电池}应怎样放置才能共同推动？画出\omterm{def:g3:first-electric-circuit:battery}{电池}队列，
并在每节\omterm{def:g3:first-electric-circuit:battery}{电池}上标出 $+$ 和 $-$。
\end{exercise}

\begin{exercise}[$\star$]\label{exo:g4:series-circuits:7}
沿手电筒的回路追踪：说出电流经过的各个元件，
从\omterm{def:g3:first-electric-circuit:battery}{电池}的 $+$ 极出发，回到它的 $-$ 极。
\end{exercise}

\begin{exercise}[$\star$]\label{exo:g4:series-circuits:8}
门铃电路中，\omterm{def:g3:first-electric-circuit:battery}{电池}、按钮和铃\omterm{def:g4:series-circuits:series}{串联}在一起。
在哪里可以再加一个按钮，使铃只有在
\emph{两个}按钮同时按下时才响？
\end{exercise}

\begin{exercise}[$\star$]\label{exo:g4:series-circuits:9}
一串十只\omterm{def:g4:series-circuits:series}{串联}的彩灯不亮了。\omterm{def:g3:first-electric-circuit:bulb}{灯泡}看上去都
完好。请描述一个有耐心的方案，用一只确定完好的\omterm{def:g3:first-electric-circuit:bulb}{灯泡}
去找出坏掉的那一只。
\end{exercise}

\begin{exercise}[$\star\star$]\label{exo:g4:series-circuits:10}
汤姆的手电筒用两节\omterm{def:g3:first-electric-circuit:battery}{电池}，发光很亮。在
黑暗中更换\omterm{def:g3:first-electric-circuit:battery}{电池}后，手电筒却不亮了——可两节新\omterm{def:g3:first-electric-circuit:battery}{电池}都是满电的，
\omterm{def:g3:first-electric-circuit:bulb}{灯泡}也完好，导线也没有松动。汤姆做了什么，
手电筒为什么仍然不亮？（见
\cref{ex:g4:series-circuits:batteries}）
\end{exercise}

\begin{exercise}[$\star\star$]\label{exo:g4:series-circuits:11}
一位工程师必须为一台机器装一个紧急停止\omterm{ex:g3:first-electric-circuit:switch}{开关}，
这台机器的电动机与其\omterm{def:g3:first-electric-circuit:battery}{电池}接在同一个单回路上。她说：
“回路上任何位置都可以。”她说得对吗？
如果这台机器的电路有支路，她的方案为什么会失败？（下
学年的章节会给它们起一个名字。）
\end{exercise}
